This paper introduced DQNGuard, a budgeted open-set decision layer for RF preliminary-action recognition from OTA captures. DQNGuard uses predicted-class calibration, variance/energy-style guard evidence, and known-only thresholding to decide whether each RF window should be retained as known PA evidence or routed as an unknown behavior.

Across WiFi, Bluetooth, and Zigbee OTA captures, DQNGuard improves the fixed-budget operating point relative to VarMax and a DQN-IDS-style confidence head. The Target--Surrogate Matrix further shows that surrogate-open calibration is target-dependent: withheld surrogate classes are not interchangeable proxies for future unknown behaviors.

These results support RF open-set recognition as a sensing and triage layer for QR-CWoS rather than as a complete response system. Future work will focus on principled surrogate selection, pooled calibration, larger RF behavior taxonomies, and downstream attack-chain integration.
